\section{符号说明}

主要符号列于表\ref{tab:符号说明}。问题二拟合时以十亿参数和十亿 tokens 计量规模，问题三计算成本时转换为参数个数和 token 数。

\begingroup
\captionsetup[table]{font={bundledsong,minusfour,bf}}
\setlength{\tabcolsep}{3pt}
\renewcommand{\arraystretch}{1.10}
\begin{longtable}{@{}>{\centering\arraybackslash}p{0.28\textwidth}>{\raggedright\arraybackslash}p{0.69\textwidth}@{}}
\caption{主要符号及含义}\label{tab:符号说明}\\
\toprule
符号 & 含义\\
\midrule
\endfirsthead
\caption{主要符号及含义（续）}\\
\toprule
符号 & 含义\\
\midrule
\endhead
\bottomrule
\endfoot
$N$ & 模型包含的参数数量\\
$D$ & 用于模型训练的数据总量\\
$L,E$ & 分别表示验证交叉熵损失与不可约损失\\
$Q$ & 数据质量评分，分值越高则质量越好\\
$Q_0$ & 质量基线；问题二说明来源间的参考水平，问题三采用参考配比加权质量\\
$\bm p$ & 各领域训练占比构成的向量，各项非负且总和为 1\\
$A,B$ & 参数规模项、数据规模项的标度系数\\
$\alpha,\beta$ & 参数规模项、数据规模项的幂指数\\
$C,C_q$ & 质量缺口系数，问题三记为 $C_q$\\
$\gamma$ & 质量缺口的幂指数\\
$\varepsilon_N,\varepsilon_D,\varepsilon_Q$ & 损失对参数量、数据量和质量的弹性\\
$M(\bm p)$ & 相对参考配比的损失平移修正\\
$C_{\text{train}},C_Q,C_{\text{attn}}$ & 基础训练、质量提升和注意力开销\\
$C_{\text{budget}},C_{\mathrm b}$ & 总算力预算，问题三简记为 $C_{\mathrm b}$\\
$g(Q)$ & 单个 token 的质量处理成本\\
$\eta$ & 注意力开销系数；参数个数作为计数代入成本公式\\
$L_{\text{ctx}},L_{\text{crit}}$ & 上下文长度、注意力与训练开销相等时的长度\\
$\lambda(L_{\text{ctx}})$ & 注意力开销与基础训练开销的比值\\
$\tau(Q)$ & 质量处理成本折算的等效附加参数规模\\
$\Theta(C_{\mathrm b})$ & 质量提升完成度\\
$F(t),K$ & 能力前沿及 logistic 饱和上界\\
$t,t_0$ & 以 2019 年为原点的时间、logistic 拐点时间\\
$r$ & logistic 增长率\\
$b_P$ & 参数量每扩大十倍对应的回归得分变化\\
$b_t$ & 面板回归的年度得分变化系数\\
$\alpha_t$ & 面板回归的年份固定效应\\
$\bar B$ & Leaderboard 六维平均得分\\
$Q_i,Q_d,\bar Q(\bm p)$ & 依次表示样本评分、领域平均评分和按训练配比加权的质量指数\\
$\bm P,\bm L$ & 分别由各组训练配比及其对应的验证损失构成的矩阵\\
$\bm A,\bm b$ & 配比回归中的系数矩阵和截距向量\\
$\lambda$ & 配比回归的岭惩罚系数，设为 $10^{-3}$\\
$\tau$ & 判断两个指标的标准化得分是否冲突的阈值，设为 $1.5$\\
$w_j,e_j$ & 指标 $j$ 对应的熵权系数与信息熵值\\
$W$ & Kendall 协同系数，用于衡量指标排序的一致程度\\
$\rho_c,\rho_0$ & 指标评价的冲突比例及指标独立条件下的参考冲突比例\\
$R^2$ & 决定系数\\
RMSE、MAE & 均方根误差、平均绝对误差\\
MAPE & 平均绝对百分比误差\\
\end{longtable}
\endgroup
